%% Appendix C: Code Listings: appendices/appendixC.tex

% ==============================================================================
% CONTENT BLUEPRINT — Appendix C: Source Code Listings
% ==============================================================================
% Full source of every analysis / processing script referenced in
% Chapters 4 and 5.
%
% Recommended content:
%   • One subsection per script. Subsections titled by the script's name.
%   • Each listing:
%       – Choice of language style is whichever matches the script:
%         style=pythonstyle, style=fortranstyle, style=matlabstyle, ...,
%         or lstinputlisting{filename} if the file is versioned in the
%         accompanying Git repository.
%       – Caption above the listing, with version hash or tag if relevant.
%       – Label is lst:<script>_<purpose> for easy cross-reference.
%       – Inline comments inside the code MUST explain non-obvious lines.
%   • State, somewhere, the canonical public repository (GitHub + Zenodo
%     DOI) where the as-shipped code lives. Reference this once near the
%     start of Appendix C.
%   • Cite any third-party library used (NumPy, SciPy, pymatgen, ASE,
%     Quantum ESPRESSO, VASP, ...)
%     — full citation goes in the bibliography, but a quick pointer
%       inside the appendix helps reviewers find dependencies.
%   • If you write Fortran, MATLAB, or Python together, choose the
%     listings = flexible approach so the listings indent sensibly with
%     double-spacing of the rest of the thesis (see preamble/packages.tex
%     for the per-language style definitions; all use \footnotesize with
%     \setstretch{1.0} so the lines stay single-spaced inside the listing
%     even when the chapter is double-spaced).
%
% Conventions:
%   • Keep listing bodies NO WIDER than 0.9 \textwidth.
%   • For long scripts, prefer to show only the function being described
%     and reference the full file in the Git repo.
%   • Do NOT embed edited screenshots of IDE windows. Listings are
%     text-only and version-controlled.
% ==============================================================================

\chapter{Source Code Listings}
\label{app:code}

\section{Fortran 95 Numerical Integration Script}
Listing~\ref{lst:fortran_code} shows a sample Fortran 95 subroutine implementing the trapezoidal rule for numerical integration.

\begin{lstlisting}[style=fortranstyle, caption={Fortran 95 implementation of the trapezoidal rule.}, label={lst:fortran_code}]
subroutine trapezoidal(f, a, b, n, result)
  implicit none
  ! Dummy arguments
  real, external :: f
  real, intent(in) :: a, b
  integer, intent(in) :: n
  real, intent(out) :: result
  
  ! Local variables
  real :: h, x, sum_val
  integer :: i
  
  h = (b - a) / float(n)
  sum_val = 0.5 * (f(a) + f(b))
  
  do i = 1, n-1
     x = a + float(i) * h
     sum_val = sum_val + f(x)
  end do
  
  result = sum_val * h
end subroutine trapezoidal
\end{lstlisting}

\section{MATLAB Data Plotting Script}
Listing~\ref{lst:matlab_code} demonstrates a MATLAB script to load calculation results and generate a publication-quality figure.

\begin{lstlisting}[style=matlabstyle, caption={MATLAB script to plot strained bandgaps.}, label={lst:matlab_code}]
% MATLAB script to plot bandgap vs strain
data = load('strain_bandgap.txt');
strain = data(:, 1);
bandgap = data(:, 2);

figure('Position', [100, 100, 600, 400]);
plot(strain, bandgap, '-or', 'LineWidth', 2, 'MarkerSize', 8);
grid on;
xlabel('Biaxial Strain (%)', 'FontSize', 12, 'FontWeight', 'bold');
ylabel('Bandgap Energy (eV)', 'FontSize', 12, 'FontWeight', 'bold');
title('Bandgap Tuning of Monolayer MoS_2', 'FontSize', 14);
saveas(gcf, 'bandgap_plot.png');
\end{lstlisting}
